\section{Direct Closure and the Four-Body Diagnostic}
\label{sec:positive-forward-four-body-diagnostic}

\subsection{Where the Factorial Tower Enters the Four-Body Circuit}
\label{subsec:factorial-tower-four-body-entry}

The Stirling/binomial route refines the first body of the four-body circuit.  The raw mixed-derivative expansion creates coefficients
\[
  {m\choose a}{\alpha\choose\beta},
\]
while the factorial tower weight
\[
  w_{m,\alpha}(\tau,\rho)={\tau^m\rho^{|\alpha|}\over m!\,\alpha!}
\]
absorbs those coefficients exactly.  The first body therefore receives an exact tower-normalization payment rather than an informal combinatorial estimate.

With moving time and space radii, the same tower produces the signed exchange
\[
  dL_S+dD_S^{\rm rad}=dJ_{SQ},
  \qquad D_S^{\rm rad}\ge0.
\]
This is the coupled analytic-radius spend.  It is the amount of tower radius paid by the selected scale channel as the four-body circuit moves toward a terminal endpoint.

The diagnostic question passed to the next body is the same-carrier reserve question.  The selected activity \(a_j(s)\) has to be priced by a same-carrier action \(B_j(s)\),
\[
  a_j(s)\le \lambda B_j(s)+\ell_j(s),
  \qquad \int_{-1}^{0}\ell_j(s)\,ds=o_j(1),
\]
and the same action has to satisfy the terminal-family reserve
\[
  \lim_{N\to\infty}\sum_{j:\,r_j\le2^{-N}}\int_{-1}^{0}B_j(s)\,ds=0.
\]
This is the current forward-gold wall named by the active thread: coupled radius spend plus same-carrier reserve.  The four-body circuit therefore routes the factorial tower into a precise proof burden instead of leaving Stirling/binomial cancellation as a side calculation.

The direct proof target for Navier--Stokes is exact. One wants estimates that
start from the original smooth divergence-free datum, remain on the same
Navier--Stokes solution, and satisfy a continuation criterion for the classical
solution beyond any proposed finite terminal time. In its cleanest form, the
program has a simple rhythm: identify the dangerous channel, pay for it with an
equation-level estimate, recover derivative control, and continue the solution.

The four-body circuit was one serious way of organizing that positive-forward
program. It is used here as a diagnostic exercise. Its purpose is to show
what happens when a forward smoothness proof is pursued as far as the available
structure can push it. The circuit is useful because each of its bodies is a real
mathematical demand, and each demand exposes the same underlying obstruction
from a different angle.

The four bodies are these.

\begin{enumerate}[leftmargin=*]
\item \textbf{Tail control.} High-frequency transfer must become summable, or
      the positive transfer tail must be dominated by the viscous tail on the
      same selected frequencies.
\item \textbf{Energy/enstrophy control.} The known energy law must be lifted
      into derivative-level control satisfying a continuation criterion.
\item \textbf{Exact-object capture.} Controlled approximations have to produce
      a same-solution terminal object rather than a detached limit object.
\item \textbf{Derivative recovery.} The captured object must return the local
      or global derivative control required by the continuation criterion.
\end{enumerate}

In the tail item, domination means a proved inequality for the selected
high-frequency tail, not an active-scale picture. For spectral balance variables
\(E(k,t)\) and \(T_+(k,t)=\max(T(k,t),0)\), the exact object is
\[
  \mathcal D_K(t)
  :=
  2\nu\int_K^\infty k^4E(k,t)\,dk
  -
  \int_K^\infty k^2T_+(k,t)\,dk,
\]
and the proof must supply \(\mathcal D_K(t)\ge0\), or an integrable bound on
its negative part, on the same tail used by the later compactness or
continuation step.

The circuit has four genuine theorem inputs.
High-frequency damping is built into viscosity. Energy is genuinely controlled.
Compactness can extract a terminal object only after the extraction topology,
low-mode compactness, high-frequency tail smallness, and same-solution
attachment have been supplied. Derivative recovery must be a named continuation
readout on that same surface. The problem appears when the four payments are
placed in a loop.

Tail control asks for a transfer estimate whose conclusion dominates nonlinear
cascade. The energy identity supplies
\[
  \sup_{0\le t<T}\|u(t)\|_{L^2}^2
  +2\nu\int_0^T\|\nabla u(t)\|_{L^2}^2\,dt
  \le \|u_0\|_{L^2}^2,
\]
which is a real payment, yet it pays an integrated first-derivative bill. It does
not by itself pay the pointwise high-frequency transfer bill needed for a full
scale barrier.

The enstrophy lift exposes the same issue in a sharper form. Set
\[
  Y(t)=\frac12\|\nabla u(t)\|_{L^2}^2.
\]
Differentiating the equation, testing the differentiated equation against
\(-\Delta u\), integrating by parts, and using \(\nabla\cdot u=0\) gives
\[
  \frac12\frac{d}{dt}\|\nabla u(t)\|_{L^2}^2
  +\nu\|\Delta u(t)\|_{L^2}^2
  =\int_{\mathbb T^3}(u\cdot\nabla)u\cdot\Delta u\,dx.
\]
H\"older, Sobolev, interpolation, and Young give the conditional differential
inequality
\[
  Y'(t)\le C_\nu Y(t)^3.
\]
This inequality gives the local continuation mechanism, and it also shows the unpaid bill:
the vortex-stretching source can still feed the derivative quantity that
viscosity is trying to damp.

Compactness then asks for the exact same-solution object produced by this
partially controlled evolution. The low/high split gives the right idea: low
modes become finite-dimensional after time compactness, and high modes become
small under a spectral tail estimate. The dependency order is the obstruction.
Compactness consumes the tail and source controls that the earlier bodies were
trying to prove.

Derivative recovery closes the loop. Once an exact object is captured, one wants
Sobolev, local-critical, endpoint, geometric, or packet-level control to recover
classical continuation. Each recovery mechanism must stay on the same
Navier--Stokes surface. A recovery mechanism that changes the surface, imports a
geometric drain from a different carrier, or detaches the terminal object from
the original solution only explains an obstruction until it is reattached to the
original solution and proves direct continuation.

The governing route of this paper is the selected-record class test.  The
four-body circuit contributes the obstruction shape created by positive-forward
closure attempts:
\[
  \text{tail/source/exact-object/recovery demand}
  \quad\longrightarrow\quad
  \text{persistent terminal survivor}.
\]
The survivor changes names across branches: high-frequency tail, source reserve,
donor residue, lower-prefix memory, packet defect, pressure-sustain object,
endpoint remnant, or export boundary. The pattern is the same. A forward proof
tries to delete the survivor. The survivor returns as a terminal object whose
proof role must be classified.

This is the point where the paper changes proof polarity. The next question is
no longer only how to erase every survivor by a stronger positive estimate. The
next question is what the survivor is allowed to be as a terminal witness for the
same Navier--Stokes solution. That question leads to the class-membership exit
method: a terminal obstruction must be tested as a same-solution witness, and its
failure must be read through the continuation requirements needed for membership.

The four-body diagnostic supplies the first branch in the larger
manuscript pattern. Each later branch will be treated in the same way: state the
positive aim, carry out the calculation, identify the terminal survivor, locate
the point where positive-forward closure stalls, and then classify the survivor
inside the CM-exit program.
